\section{Why Tabular Data Breaks Text-Centric Agentic RAG}
\label{sec:diagnosis}

This section diagnoses the failure mechanics produced by treating tabular sources as if they were passages. Each mechanism we identify here is addressed by a specific MOMO mechanism in Section~\ref{sec:architecture}; the closing paragraph of this section ties each failure mode to its intervention.

\paragraph{Multi-turn source commitment.} A useful table is rarely usable in one turn. The agent must peek the schema, choose columns, write a query, recover from a parsing error, validate the result against the question, and decide whether the source is exhausted. Standard agentic-RAG loops~\cite{yao2023react,trivedi2023ircot} have no primitive for ``I have committed to this source --- finish or release it.'' Without that primitive, the agent treats every turn as if the source were new, which manifests as either over-investment (continuing to query a source that will not produce evidence) or under-investment (switching sources before extracting available evidence).

\paragraph{Stopping-rule absence.} Without an explicit commitment contract, the agent oscillates: re-search $\to$ re-peek $\to$ re-query $\to$ switch source $\to$ re-search. This loop is not a planning failure --- the agent often has the correct candidate dataset --- it is a stopping-rule failure. The agent does not know when retrieving more sources stops being useful and execution should take over. In our traces, this is the single largest source of timeouts: the budget runs out mid-search.

\paragraph{Context pollution.} Tabular wrangling artifacts are large. A typical \texttt{peek\_file} response includes the schema (tens to hundreds of columns with names and types), a snippet of rows, optional column statistics, and an LLM-generated description. A failed SQL query returns the query text and a traceback that may include the offending row. Across a few candidate sources, the planner's context fills with these artifacts. Once the context is dominated by raw data, the LLM's reasoning quality on the original question degrades, which in turn produces worse query choices, which produce more failures, which produce more artifacts. The system enters a context-pollution spiral.

\paragraph{File-handle errors.} The tool surface accepts either a logical \texttt{dataset\_id} or a concrete \texttt{s3\_uri}/path. A planner that confuses the two loses turns to ``file not found'' errors. These are not reasoning failures; they are bookkeeping failures. In our traces, a non-trivial fraction of total wasted budget is spent on retrying tool calls that referenced a dataset by ID without the exact file path. The agent eventually figures it out, but each instance costs a turn and adds an unhelpful traceback to the context.

\paragraph{Source drift.} When peeking is cheap, planners drift into adjacent datasets that look tangentially relevant. Without scoping discipline, the model rationalizes each detour --- ``let me check this other table to confirm'' --- and the rationalizations accumulate. By the time the agent returns to its primary source, several turns have been spent on data unrelated to the question.

\paragraph{Premature submit.} The mirror image of stopping-rule absence. With no enumerated \texttt{missing\_outputs} discipline, the planner sometimes submits an answer that covers only some of the requested outputs. Two examples: a question asking for 2019--2021 figures answered with only 2019--2020; a question asking for ``the top 5 counties'' answered with only the top 1. Nothing in the agent loop forces the planner to enumerate what it has not yet found before declaring completion.

\paragraph{From failure mechanisms to MOMO mechanisms.} Each of the six failure mechanisms above motivates a specific MOMO design choice (covered in detail in Section~\ref{sec:architecture}). \emph{Multi-turn source commitment} and \emph{stopping-rule absence} motivate the bounded inspect subagent: an explicit per-source budget makes commitment a first-class concept, and a budget-exhausted return forces the planner to acknowledge stopping rather than oscillate. \emph{Context pollution} motivates the context firewall: raw schemas and tracebacks live inside the worker and never enter the planner's conversation. \emph{File-handle errors} motivate \texttt{\_FileReferenceGuard}, a steering handler that blocks tool calls referencing a dataset by ID without the exact \texttt{s3\_uri}, surfacing bookkeeping errors as immediate contract violations rather than silent retries. \emph{Source drift} motivates \texttt{\_InspectSourceGuard}, which restricts an inspect worker to the \texttt{source\_family\_ids} declared in its contract. \emph{Premature submit} motivates the partial-success protocol, in which a worker that has not satisfied every \texttt{required\_output} must return \texttt{status="partial"} with an explicit \texttt{missing\_outputs} list, forcing the planner to either retry or accept the gap explicitly before answering.
